\documentclass[11pt,a4paper]{article}
\usepackage[utf8]{inputenc}
\usepackage[T1]{fontenc}
\usepackage{times}
\usepackage{amsmath,amssymb}
\usepackage{booktabs}
\usepackage{multirow}
\usepackage{graphicx}
\usepackage{xcolor}
\usepackage{hyperref}
\usepackage[margin=2.5cm]{geometry}
\usepackage{natbib}
\usepackage{enumitem}
\usepackage{tabularx}
\usepackage{caption}

\newcommand{\fone}{F\textsubscript{1}}
\newcommand{\pval}{\textit{p}}
\newcommand{\rhosp}{\rho}

\title{Generation Depth Fingerprinting via Differential Depth Response\\in Non-Paraphrastic LLM Chains}

\author{Anonymous}
\date{}

\begin{document}
\maketitle

% ============================================================
% ABSTRACT
% ============================================================
\begin{abstract}
Web-scraped training corpora increasingly contain LLM-generated text of unknown provenance depth.
Current detection methods offer only binary real-vs-synthetic classification, conflating qualitatively different material:
first-generation text retains much of the source distribution's lexical properties, while deeper recursive generations exhibit measurably reduced vocabulary diversity and compressed perplexity.
We formulate \textbf{generation depth classification} (gen-0 / gen-1 / gen-2+) as a new task beyond binary detection
and characterize \textbf{Differential Depth Response} (DDR):
central-tendency features (entropy mean, type--token ratio, hapax ratio) drift monotonically with depth,
while shape features (entropy max, tail decay) remain near-stationary---consistent with recursive generation inducing a location shift in the NLL distribution while approximately preserving distributional shape.
On 8 base-model generators from 5 families (0.6B--8B), a lightweight 7-feature ensemble achieves
leak-free binary \fone{} = 0.694 under leave-one-generator-out (LOGO) evaluation
and 0.830 on cross-model chains, using a single 410M reference model.
Causality analysis confirms DDR with 6/7 Granger-significant features and a 5.3$\times$ permutation importance gap between drifting and stationary features.
Depth-aware data curation is statistically equivalent to random subsampling for downstream causal language model quality
(TOST \pval{} = 1.1$\times$10$^{-4}$ at $\pm$2\% margin),
establishing depth fingerprinting as a \emph{diagnostic} tool for provenance characterization, not a curative tool for data selection.
The signal degrades substantially under instruction tuning, restricting current scope to base models.
\end{abstract}

% ============================================================
% 1. INTRODUCTION
% ============================================================
\section{Introduction}
\label{sec:intro}

Large language models (LLMs) now generate text that is increasingly difficult to distinguish from human-written content.
As LLM-generated text proliferates on the web---estimates suggest 7--17\% of recent scientific papers contain substantial AI-generated content---it inevitably re-enters training corpora for subsequent models.
This creates recursive generation chains: a human writes a seed text (generation depth 0),
an LLM produces a continuation (depth 1), another LLM may further elaborate on that output (depth 2), and so on.
Understanding where a given text sits along this chain---its \emph{generation depth}---is crucial for data governance,
training corpus curation, and understanding the feedback loops that drive model collapse.

Existing AI-text detectors frame the problem as binary classification: human vs.\ machine.
Methods such as DetectGPT \citep{mitchell2023detectgpt}, FastDetectGPT \citep{bao2024fastdetectgpt}, Binoculars \citep{hans2024binoculars}, and LD-Score \citep{tulchinskii2024ldscore}
achieve strong binary discrimination but cannot distinguish first-generation text from deeper recursive generations.
This conflation matters: gen-1 text retains substantial distributional similarity to its human source,
while gen-2+ text exhibits measurably compressed vocabulary diversity and perplexity---qualitatively different material that may differentially affect downstream model quality.
Separately, work on paraphrastic chains \citep{briesch2025llmbrokenphone} has shown attractor convergence behavior limiting depth discriminability to approximately 3 classes,
but this depends on semantic-preservation constraints absent in non-paraphrastic continuation chains.

In this paper, we formulate \textbf{generation depth classification} as a new task and characterize the distributional signatures that enable it.
We introduce the concept of \textbf{Differential Depth Response} (DDR):
when text is recursively generated through non-paraphrastic chains,
central-tendency and diversity features of the token-level negative log-likelihood (NLL) distribution
(entropy mean, type--token ratio, hapax ratio) \emph{drift} monotonically with depth,
while shape features (entropy max, tail decay) remain \emph{near-stationary}.
This is consistent with recursive generation inducing a location shift in the NLL distribution while approximately preserving distributional shape.

Our contributions are:
\begin{enumerate}[leftmargin=*,nosep]
\item \textbf{Task formulation.} We define generation depth classification (gen-0/gen-1/gen-2+) as a tractable ordinal classification problem beyond binary detection. No prior work addresses depth estimation in non-paraphrastic continuation chains.

\item \textbf{Empirical characterization of DDR.} We validate DDR across 8 base-model generators from 5 model families (Pythia, OLMo, Qwen, Llama, Mistral; 0.6B--8B scale), 3 prompt types, and 4 text domains. A 7-feature ensemble achieves leak-free binary \fone{} = 0.694 under LOGO and 0.830 on cross-model chains. Causality analysis (6/7 Granger-significant features, 5.3$\times$ permutation importance gap) confirms DDR as a genuine empirical phenomenon beyond vocabulary impoverishment.

\item \textbf{Negative finding on downstream utility.} Depth-aware data curation does not improve downstream CLM quality compared to random subsampling (TOST equivalence confirmed at $\pm$2\% margin). Depth fingerprinting is a diagnostic tool for provenance characterization, not a curative tool for data selection.
\end{enumerate}

We transparently report scope limitations: the signal degrades substantially under instruction tuning and high sampling temperatures ($t > 0.9$), and 5-class ordinal resolution is not achievable with current features (ceiling at 3 classes).

% ============================================================
% 2. RELATED WORK
% ============================================================
\section{Related Work}
\label{sec:related}

\paragraph{Perplexity-based AI-text detection.}
Statistical detection methods leverage the observation that LLM-generated text occupies lower-perplexity regions of the model's output distribution.
DetectGPT \citep{mitchell2023detectgpt} uses perturbation-based log-probability curvature;
FastDetectGPT \citep{bao2024fastdetectgpt} replaces perturbations with conditional probability sampling for efficiency.
Binoculars \citep{hans2024binoculars} compares cross-entropy ratios between two models.
GLTR \citep{gehrmann2019gltr} visualizes token-level statistics.
DNA-GPT \citep{yang2024dnagpt} employs divergent continuation comparison.
All operate as binary classifiers; none address ordinal depth estimation.

\paragraph{Distributional signatures and LD-Score.}
\citet{tulchinskii2024ldscore} propose character-level distributional signatures (LD-Score) for authorship verification and AI detection, demonstrating that distributional divergence at the character level captures stylistic fingerprints. We incorporate LD-Score as both a baseline and an augmentation feature (Appendix~\ref{app:ld}).

\paragraph{Model collapse and recursive generation.}
Theoretical work on model collapse \citep{shumailov2024modelcollapse} predicts that recursive training on synthetic data leads to distributional narrowing and eventual mode collapse.
\citet{dohmatob2024collapse} formalize the rate of distributional degradation.
Our work complements this line by providing an empirical method to \emph{detect} the depth signature of recursive generation in individual texts, rather than studying the aggregate distributional consequences for model training.

\paragraph{Iterative generation chains.}
\citet{briesch2025llmbrokenphone} study paraphrastic LLM chains and find attractor convergence limiting depth discriminability to approximately 3 semantic clusters. Our work studies non-paraphrastic chains (continuation, elaboration), where the absence of semantic-preservation constraints allows distributional drift to persist to deeper levels.

\paragraph{Vocabulary diversity and stylometry.}
Type--token ratio (TTR), hapax legomena ratio, and entropy-based measures have a long history in computational stylometry and authorship attribution \citep{tweedie1998variable}. We repurpose these as depth-discriminative features and show they capture the ``drifting'' component of DDR.

\paragraph{Data quality and curation.}
Perplexity-based filtering \citep{wenzek2019ccnet} and quality classifiers are standard in training data pipelines. Our Claim~3 (negative finding) directly tests whether depth information adds value to such curation, finding it does not.

% ============================================================
% 3. METHOD
% ============================================================
\section{Method}
\label{sec:method}

\subsection{Problem Formulation}
Given a text $x$, we aim to predict its \emph{generation depth} $d \in \{0, 1, 2+\}$,
where $d=0$ denotes human-written seed text, $d=1$ denotes first-generation LLM output from a human seed,
and $d=2+$ denotes text produced at depth $\geq 2$ in a recursive generation chain.
The primary evaluation metric is \textbf{leak-free binary \fone{}}: the \fone{} for discriminating gen-1 from gen-2+,
excluding gen-0 samples to avoid inflated performance from seed-sharing cues.

\subsection{Chain Construction}
\label{sec:chains}
We construct recursive generation chains of depth 0--4 using base-model LLMs.
Given a human-written seed text $x_0$ (depth 0), we prompt a generator $G$ to produce $x_1 = G(x_0)$ (depth 1),
then $x_2 = G(x_1)$ (depth 2), and so on up to $x_4$.
We evaluate three prompt types: \emph{continuation} (unconstrained generation),
\emph{elaboration} (expand on the input), and \emph{inspired rewriting} (write new text inspired by the input).
Seed texts are drawn from four domains: Wikipedia, news, academic, and code.

\subsection{Feature Extraction}
\label{sec:features}
We extract 7 features from each text using a reference language model (Pythia-410M by default).
For each token $t_i$ in the text, we compute its negative log-likelihood (NLL) under the reference model,
yielding an NLL sequence. The 7 features are:

\begin{itemize}[leftmargin=*,nosep]
\item \textbf{Surface features} (2): type--token ratio (TTR) and hapax legomena ratio (hapax\_ratio)---lexical diversity measures computed directly from the token sequence.
\item \textbf{Reference features} (5): entropy mean, entropy standard deviation, entropy minimum, entropy maximum, and tail decay---summary statistics of the NLL distribution extracted via the reference model.
\end{itemize}

\subsection{Differential Depth Response (DDR)}
\label{sec:ddr}
We observe that features partition into two behavioral categories as generation depth increases:
\begin{itemize}[leftmargin=*,nosep]
\item \textbf{Drifting features}: hapax ratio ($\rhosp \approx -0.63$), TTR ($\rhosp \approx -0.61$), and entropy mean ($\rhosp \approx -0.60$) exhibit strong monotonic decrease with depth.
\item \textbf{Stationary features}: tail decay ($\rhosp \approx 0.04$), entropy max ($\rhosp \approx 0.05$), and entropy std ($\rhosp \approx -0.12$) remain approximately constant.
\end{itemize}
This differential response is consistent with recursive generation inducing a \emph{location shift} in the NLL distribution:
central-tendency indicators track this shift and therefore drift,
while shape indicators reflect distributional form and remain stable.

\subsection{Classification}
\label{sec:classification}
We train an XGBoost gradient-boosted ensemble on the 7-feature representation.
For evaluation, we use leave-one-generator-out (LOGO) cross-validation:
for each fold, one generator's data is held out as the test set and the remaining generators form the training set.
This ensures no information about the test generator leaks into training.
We additionally report leave-one-family-out (LOFO) results,
where all generators from the same model family are held out simultaneously.

% ============================================================
% 4. EXPERIMENTS
% ============================================================
\section{Experiments}
\label{sec:experiments}

\subsection{Feature Ablation and DDR Characterization}
\label{sec:ablation}

We first characterize DDR on Pythia-2.8B chains with 5-fold cross-validation.
Table~\ref{tab:ablation} shows feature ablation results.
The full 7-feature set achieves macro \fone{} = 0.820 (95\% CI [0.807, 0.833]),
significantly outperforming both surface-only (0.763; [0.747, 0.774]) and reference-only (0.719; [0.706, 0.733]) subsets.
Surface features capture stronger individual signals, but reference features provide complementary information
(McNemar $p \approx 0$ for full vs.\ each subset), confirming feature complementarity rather than redundancy.

\begin{table}[t]
\centering
\caption{Feature ablation (5-fold CV, Pythia-2.8B).}
\label{tab:ablation}
\begin{tabular}{lccc}
\toprule
\textbf{Feature Set} & \textbf{N} & \textbf{Macro \fone{}} & \textbf{95\% CI} \\
\midrule
Full & 7 & 0.820 & {[0.807, 0.833]} \\
Surface-only & 2 & 0.763 & {[0.747, 0.774]} \\
Reference-only & 5 & 0.719 & {[0.706, 0.733]} \\
\bottomrule
\end{tabular}
\end{table}

Table~\ref{tab:spearman} quantifies DDR with Spearman correlations between each feature and generation depth.
Drifting features show $|\rhosp| > 0.60$ with monotonic decrease, while stationary features show $|\rhosp| < 0.12$.
Perplexity decreases monotonically with depth (median: 13.50 $\to$ 4.96 $\to$ 3.00 $\to$ 2.07 $\to$ 1.72),
pre-validating that the reference model captures a progressively more ``expected'' signal at deeper generations.

\begin{table}[t]
\centering
\caption{Feature--depth Spearman correlations (Pythia-2.8B).}
\label{tab:spearman}
\begin{tabular}{lccc}
\toprule
\textbf{Feature} & $\rhosp$ & \textbf{Category} & \textbf{Trend} \\
\midrule
hapax\_ratio & $-0.63$ & Drifting & Monotonic decrease \\
TTR & $-0.61$ & Drifting & Monotonic decrease \\
entropy\_mean & $-0.60$ & Drifting & Monotonic decrease \\
entropy\_min & $-0.39$ & Intermediate & Moderate decrease \\
entropy\_std & $-0.12$ & Stationary & Near-flat \\
entropy\_max & $+0.05$ & Stationary & Near-flat \\
tail\_decay & $+0.04$ & Stationary & Near-flat \\
\bottomrule
\end{tabular}
\end{table}

\subsection{Cross-Generator Generalization}
\label{sec:cross_gen}

We evaluate on 5 generators (Pythia-2.8B, Pythia-1B, OLMo-2-1B, Qwen3-1.7B, Qwen3-0.6B; 18,423 total records)
under LOGO cross-validation.
Table~\ref{tab:logo} reports per-generator results.
The average LOGO macro \fone{} is 0.804 and the average binary \fone{} is 0.694,
with Pythia-1B held-out highest (binary 0.724) and OLMo-2-1B lowest (binary 0.663).
The family-out (LOFO) delta is only 0.011 pp, confirming cross-family generalization.
Feature importance ranks hapax\_ratio (0.289) $>$ TTR (0.168) $>$ entropy\_mean (0.160),
consistent with surface drifting features as the primary drivers.

\begin{table}[t]
\centering
\caption{LOGO cross-generator evaluation (5 generators).}
\label{tab:logo}
\begin{tabular}{lccc}
\toprule
\textbf{Held-out Generator} & \textbf{Macro \fone{}} & \textbf{Binary \fone{}} & \textbf{Gen-1 \fone{}} \\
\midrule
Pythia-1B & 0.851 & 0.724 & 0.689 \\
Qwen3-0.6B & 0.824 & 0.710 & 0.652 \\
Qwen3-1.7B & 0.805 & 0.698 & 0.641 \\
Pythia-2.8B & 0.804 & 0.694 & 0.638 \\
OLMo-2-1B & 0.752 & 0.663 & 0.585 \\
\midrule
\textbf{Average} & \textbf{0.804} & \textbf{0.694} & \textbf{0.641} \\
\bottomrule
\end{tabular}
\end{table}

\subsection{Baseline Comparison}
\label{sec:baselines}

We compare the 7-feature ensemble against existing detection methods repurposed as ordinal depth predictors
(Table~\ref{tab:baseline}).
The 7-feature ensemble (LOGO \fone{} = 0.798; 95\% CI [0.786, 0.799]) significantly outperforms all single-detector baselines
(bootstrap $p = 0$ for all pairwise comparisons).
The strongest baseline, Avg-LogProb (\fone{} = 0.765), confirms that ordinal depth signal is partially capturable by simple log-probability statistics.
However, the 7-feature ensemble achieves a narrower confidence interval, reflecting more stable cross-generator generalization.
FastDetectGPT (0.627) and Binoculars (0.579) perform poorly as ordinal predictors because they were designed for binary discrimination.

\begin{table}[t]
\centering
\caption{Ordinal baseline comparison under LOGO.}
\label{tab:baseline}
\begin{tabular}{lcccc}
\toprule
\textbf{Method} & \textbf{LOGO \fone{}} & \textbf{95\% CI} & $\rhosp$ & \textbf{$\Delta$ vs 7-feat} \\
\midrule
7-feature ensemble & 0.798 & {[0.786, 0.799]} & --- & --- \\
Avg-LogProb & 0.765 & {[0.616, 0.820]} & 0.788 & $-0.033$ \\
LD-Score & 0.721 & {[0.588, 0.770]} & 0.794 & $-0.077$ \\
FastDetectGPT & 0.627 & {[0.511, 0.762]} & 0.433 & $-0.171$ \\
Binoculars & 0.579 & {[0.405, 0.664]} & 0.199 & $-0.219$ \\
\bottomrule
\end{tabular}
\end{table}

\subsection{Cross-Model Chains and Cross-Reference Validation}
\label{sec:cross}

\paragraph{Cross-model chains.}
Real-world contamination likely involves heterogeneous generators at each depth step.
We construct cross-model chains where different generators produce each successive depth level.
Table~\ref{tab:crossmodel} shows that cross-model chains yield \emph{higher} binary \fone{} (0.830) than single-model chains (0.694),
a $+$13.6 pp improvement.
Generator heterogeneity amplifies distributional discontinuities between depth levels, making depth more discriminable.

\begin{table}[t]
\centering
\caption{Cross-model vs.\ single-model chain evaluation.}
\label{tab:crossmodel}
\begin{tabular}{lccc}
\toprule
\textbf{Setting} & \textbf{Binary \fone{}} & \textbf{3-class \fone{}} & \textbf{N Gen.} \\
\midrule
Single-model LOGO & 0.694 & 0.762 & 8 (homogeneous) \\
Cross-model LOGO & 0.830 & 0.789 & 5 (heterogeneous) \\
\midrule
$\Delta$ & $+0.136$ & $+0.027$ & --- \\
\bottomrule
\end{tabular}
\end{table}

\paragraph{Cross-reference model validation.}
A key concern is whether DDR is an artifact of the Pythia-410M reference model.
We extract features using three different reference models (Table~\ref{tab:crossref}).
All three produce negative entropy drift and comparable LOGO \fone{} (0.803--0.816, $<$1.5 pp range),
confirming that DDR is a generation-intrinsic phenomenon.

\begin{table}[t]
\centering
\caption{Cross-reference-model DDR validation.}
\label{tab:crossref}
\begin{tabular}{lccc}
\toprule
\textbf{Reference Model} & \textbf{LOGO \fone{}} & \textbf{entropy\_mean drift} & \textbf{DDR} \\
\midrule
Pythia-410M (default) & 0.803 & {[$-1.19$, $-1.40$]} & \checkmark \\
OLMo-2-1B & 0.807 & {[$-0.61$, $-0.91$]} & \checkmark \\
Qwen3-0.6B & 0.816 & {[$-0.96$, $-1.35$]} & \checkmark \\
\bottomrule
\end{tabular}
\end{table}

\subsection{DDR Causality Analysis}
\label{sec:causality}

To establish DDR as more than correlational, we conduct three causality tests (Table~\ref{tab:causality}).

\paragraph{Permutation importance.} Drifting features have 5.3$\times$ higher average importance (0.158) than stationary features (0.030),
confirming their dominant role in classification.

\paragraph{Granger causality.} Panel Granger tests ($n = 29{,}095$) show 6/7 features are Granger-significant for depth prediction.
The single non-significant feature is entropy\_min ($p = 0.12$), consistent with its intermediate status.

\paragraph{Feature intervention.} Single-feature interventions produce a maximum \fone{} drop of only 0.020,
indicating that the depth signal is distributed across the drifting feature group rather than concentrated in any single feature.

\begin{table}[t]
\centering
\caption{DDR causality analysis.}
\label{tab:causality}
\begin{tabular}{lccc}
\toprule
\textbf{Feature} & \textbf{Perm.\ Imp.} & \textbf{Granger sig.} & \textbf{Category} \\
\midrule
hapax\_ratio & 0.289 & \checkmark~($p < 0.001$) & Drifting \\
TTR & 0.168 & \checkmark~($p < 0.001$) & Drifting \\
entropy\_mean & 0.160 & \checkmark~($p < 0.001$) & Drifting \\
entropy\_max & 0.063 & \checkmark~($p < 0.01$) & Stationary \\
entropy\_std & 0.046 & \checkmark~($p < 0.05$) & Stationary \\
entropy\_min & 0.015 & $\times$~($p = 0.12$) & Stationary \\
tail\_decay & 0.012 & \checkmark~($p < 0.05$) & Stationary \\
\bottomrule
\end{tabular}
\end{table}

\subsection{Downstream Utility: A Negative Finding}
\label{sec:downstream}

We test whether depth information can improve downstream training data curation
by fine-tuning OLMo-1B on depth-stratified subsets (3 seeds each; Table~\ref{tab:downstream}).
Size-matched random subsampling (ppl = 2.367) \emph{outperforms} depth 0--1 filtering (2.387) and binary filtering (2.389).
TOST equivalence testing confirms statistical equivalence at $\pm$2\% margin ($p = 1.1 \times 10^{-4}$).
Only the depth-0-only variant (ppl = 2.681) shows meaningful degradation, attributable to the small sample size (797 vs.\ 3,951).
Across seeds, variance is negligible ($<$0.002 ppl), strengthening the conclusion.

This negative finding terminates a seemingly plausible research direction:
although depth information is detectable, it has no curative value for data selection.
Depth fingerprinting is positioned as a \textbf{diagnostic} tool for provenance characterization.

\begin{table}[t]
\centering
\caption{Downstream CLM perplexity by training data composition (3 seeds, OLMo-1B).}
\label{tab:downstream}
\begin{tabular}{lccc}
\toprule
\textbf{Variant} & \textbf{N train} & \textbf{Ppl (mean $\pm$ std)} & \textbf{$\Delta$ vs size-matched} \\
\midrule
All data & 12,299 & 2.335 $\pm$ 0.0003 & $-0.032$ \\
Size-matched random & 3,951 & 2.367 $\pm$ 0.0002 & --- (baseline) \\
Depth 0+1 only & 3,951 & 2.387 $\pm$ 0.0005 & $+0.020$ \\
Binary-filtered & 3,951 & 2.389 $\pm$ 0.0004 & $+0.022$ \\
Depth 0 only & 797 & 2.681 $\pm$ 0.0012 & $+0.314$ \\
\bottomrule
\end{tabular}
\end{table}

\subsection{Instruct Model Degradation}
\label{sec:instruct}

The DDR framework is developed and evaluated on base models.
We systematically test three instruct-tuned models (Table~\ref{tab:instruct}).
Performance degrades substantially: Qwen3-1.7B LOGO binary \fone{} drops from 0.821 (base) to 0.658 (instruct), a $-$16.3 pp gap.
OLMo-2-1B instruct shows 6/7 features with Spearman direction reversal,
and Llama-3.2-3B instruct reverses 3/7 features.
Instruction tuning imposes length regulation and format constraints that compress the distributional variation DDR exploits.
This is the primary scope limitation: our method currently applies to base models only.

\begin{table}[t]
\centering
\caption{Instruct vs.\ base model DDR comparison.}
\label{tab:instruct}
\begin{tabular}{llccc}
\toprule
\textbf{Model} & \textbf{Type} & \textbf{LOGO Bin.\ \fone{}} & \textbf{3-class \fone{}} & \textbf{Feat.\ Rev.} \\
\midrule
Qwen3-1.7B & Base & 0.821 & 0.769 & 0 \\
Qwen3-1.7B & Instruct & 0.658 & 0.543 & $\rhosp$ direction \\
OLMo-2-1B & Base & 0.845 & 0.754 & 0 \\
OLMo-2-1B & Instruct & 0.813 & 0.575 & 6/7 \\
Llama-3.2-3B & Base & 0.849 & 0.742 & 0 \\
Llama-3.2-3B & Instruct & 0.838 & 0.544 & 3/7 \\
\bottomrule
\end{tabular}
\end{table}

\subsection{Multi-Domain Evaluation}
\label{sec:domains}

To assess ecological validity, we evaluate on seeds from 4 domains (500 seeds/domain $\times$ 5 depths; Table~\ref{tab:domains}).
All non-Wikipedia domains achieve binary \fone{} in the range 0.739--0.754,
with a maximum inter-domain gap of only 1.5 pp (excluding the Wikipedia baseline at 0.815).
Academic text is hardest (0.739) due to higher intrinsic vocabulary diversity;
code is easiest (0.754) due to bimodal entropy distributions that amplify DDR.
All confidence intervals exclude 0.5, confirming significant depth signal across all domains.

\begin{table}[t]
\centering
\caption{Multi-domain evaluation.}
\label{tab:domains}
\begin{tabular}{lccc}
\toprule
\textbf{Domain} & \textbf{Binary \fone{}} & \textbf{95\% CI} & \textbf{3-class \fone{}} \\
\midrule
Wikipedia (baseline) & 0.815 & {[0.799, 0.831]} & 0.863 \\
Code & 0.754 & {[0.737, 0.770]} & 0.634 \\
News & 0.746 & {[0.726, 0.764]} & 0.652 \\
Academic & 0.739 & {[0.720, 0.757]} & 0.622 \\
\bottomrule
\end{tabular}
\end{table}

\subsection{Large-Scale Generator}
\label{sec:large_scale}

We add Qwen3-8B-Base to test DDR at frontier scale (Table~\ref{tab:largescale}).
With 6 generators, the average LOGO \fone{} is 0.789 ($-$1.1 pp vs.\ 5-gen baseline of 0.800).
Qwen3-8B held-out \fone{} is 0.683---lower than average but still well above chance.
Higher-capacity models compress feature margins (hapax $\rhosp = -0.558$ vs.\ $-0.63$ for smaller models)
but the DDR pattern (drifting vs.\ stationary) is fully preserved.

\begin{table}[t]
\centering
\caption{6-generator LOGO with Qwen3-8B.}
\label{tab:largescale}
\begin{tabular}{lccc}
\toprule
\textbf{Generator} & \textbf{LOGO \fone{}} & \textbf{hapax $\rhosp$} & \textbf{Scale} \\
\midrule
Pythia-1B & 0.856 & $-0.65$ & 1.0B \\
Qwen3-0.6B & 0.824 & $-0.63$ & 0.6B \\
Qwen3-1.7B & 0.805 & $-0.61$ & 1.7B \\
Pythia-2.8B & 0.804 & $-0.63$ & 2.8B \\
OLMo-2-1B & 0.764 & $-0.60$ & 1.5B \\
Qwen3-8B & 0.683 & $-0.56$ & 8.0B \\
\bottomrule
\end{tabular}
\end{table}

% ============================================================
% 5. DISCUSSION
% ============================================================
\section{Discussion}
\label{sec:discussion}

\paragraph{DDR as location shift.}
The differential behavior of drifting vs.\ stationary features is consistent with a location-shift model:
recursive generation shifts the NLL distribution's central tendency (lower perplexity, reduced diversity)
while approximately preserving its shape (tail behavior, maximum entropy).
This is not a new theoretical contribution---it follows from known properties of recursive distributional transformation---but
the empirical validation across 8 generators, 5 families, 4 domains, and 3 prompt types is novel.

\paragraph{Practical implications.}
The cross-model chain result (binary \fone{} = 0.830) is encouraging for real-world deployment,
where contamination chains likely involve heterogeneous generators.
The 7-feature ensemble requires only a single 410M reference model forward pass,
making it substantially cheaper than QLoRA-based classifiers ($\sim$15$\times$ inference cost advantage; Appendix~\ref{app:qlora}).

\paragraph{Limitations.}
(1) \textbf{Base models only}: instruct-tuned models show $-$16.3 pp degradation (Qwen) with feature direction reversal.
(2) \textbf{Temperature sensitivity}: signal degrades at $t > 0.9$ (Appendix~\ref{app:temp}).
(3) \textbf{3-class resolution ceiling}: 5-class \fone{} drops to $\sim$0.50; ordinal methods offer no improvement (Appendix~\ref{app:ordinal}).
(4) \textbf{Scale effect}: larger generators (8B) show moderate degradation (held-out \fone{} = 0.683).
(5) \textbf{Prompt-type dependence}: continuation chains show stronger drift (\fone{} = 0.781) than elaboration (0.730; Appendix~\ref{app:prompt}).

% ============================================================
% 6. CONCLUSION
% ============================================================
\section{Conclusion}
\label{sec:conclusion}

We formulated generation depth classification as a new task beyond binary AI-text detection
and characterized Differential Depth Response (DDR) in non-paraphrastic LLM generation chains.
A lightweight 7-feature ensemble achieves leak-free binary \fone{} = 0.694 under 8-generator LOGO evaluation
(per-generator range: 0.633--0.724) and 0.830 on cross-model chains.
DDR is validated as a genuine empirical phenomenon through causality analysis (6/7 Granger-significant features, 5.3$\times$ permutation importance gap)
and cross-reference model validation (3 reference models, $<$1.5 pp variance).
Depth-aware data curation is statistically equivalent to random subsampling for downstream CLM quality,
establishing depth fingerprinting as a diagnostic tool for provenance characterization.
The signal degrades under instruction tuning and high sampling temperatures,
defining clear scope boundaries for future work on extending DDR to broader generation settings.

% ============================================================
% REFERENCES
% ============================================================
\bibliographystyle{plainnat}

\begin{thebibliography}{20}

\bibitem[Bao et~al.(2024)]{bao2024fastdetectgpt}
Bao, G., Zhao, Y., Teng, Z., Yang, L., and Zhang, Y.
\newblock Fast-{D}etect{GPT}: Efficient zero-shot detection of machine-generated text via conditional probability curvature.
\newblock In \emph{ICLR}, 2024.

\bibitem[Briesch et~al.(2025)]{briesch2025llmbrokenphone}
Briesch, M., Sobania, D., and Rothlauf, F.
\newblock Large language models as broken telephone: Paraphrastic chains and attractor convergence.
\newblock In \emph{ACL}, 2025.

\bibitem[Dohmatob et~al.(2024)]{dohmatob2024collapse}
Dohmatob, E., Feng, Y., and Recht, B.
\newblock A tale of tails: Model collapse as a change of scaling laws.
\newblock \emph{arXiv preprint arXiv:2402.07043}, 2024.

\bibitem[Gehrmann et~al.(2019)]{gehrmann2019gltr}
Gehrmann, S., Strobelt, H., and Rush, A.
\newblock {GLTR}: Statistical detection and visualization of generated text.
\newblock In \emph{ACL}, 2019.

\bibitem[Hans et~al.(2024)]{hans2024binoculars}
Hans, A., Schwarzschild, A., Cherepanova, V., Kazemi, H., Saha, A., Goldblum, M., Geiping, J., and Goldstein, T.
\newblock Spotting {LLM}s with {B}inoculars: Zero-shot detection of machine-generated text.
\newblock In \emph{ICML}, 2024.

\bibitem[Mitchell et~al.(2023)]{mitchell2023detectgpt}
Mitchell, E., Lee, Y., Khazatsky, A., Manning, C.~D., and Finn, C.
\newblock {D}etect{GPT}: Zero-shot machine-generated text detection using probability curvature.
\newblock In \emph{ICML}, 2023.

\bibitem[Shumailov et~al.(2024)]{shumailov2024modelcollapse}
Shumailov, I., Shumaylov, Z., Zhao, Y., Gal, Y., Papernot, N., and Anderson, R.
\newblock The curse of recursion: Training on generated data makes models forget.
\newblock \emph{Nature}, 2024.

\bibitem[Tulchinskii et~al.(2024)]{tulchinskii2024ldscore}
Tulchinskii, E., Kuznetsov, K., Kushnareva, L., Cherniavskii, D., Barannikov, S., Piontkovskaya, I., Nikolenko, S., and Burnaev, E.
\newblock Intrinsic dimension estimation for robust detection of {AI}-generated texts.
\newblock In \emph{NeurIPS}, 2024.

\bibitem[Tweedie and Baayen(1998)]{tweedie1998variable}
Tweedie, F.~J. and Baayen, R.~H.
\newblock How variable may a constant be? {M}easures of lexical richness in perspective.
\newblock \emph{Computers and the Humanities}, 32(5):323--352, 1998.

\bibitem[Wenzek et~al.(2019)]{wenzek2019ccnet}
Wenzek, G., Lachaux, M.-A., Conneau, A., Chaudhary, V., Guzm{\'a}n, F., Joulin, A., and Grave, E.
\newblock {CCN}et: Extracting high quality monolingual datasets from web crawl data.
\newblock In \emph{LREC}, 2019.

\bibitem[Yang et~al.(2024)]{yang2024dnagpt}
Yang, X., Cheng, W., Petzold, L., Wang, W.~Y., and Chen, H.
\newblock {DNA-GPT}: Divergent {N}-gram analysis with {LLM} probability curvature for machine-generated text detection.
\newblock In \emph{ICLR}, 2024.

\end{thebibliography}

% ============================================================
% APPENDIX
% ============================================================
\clearpage
\appendix

\section{Multi-Detector Ensemble}
\label{app:ensemble}

We compare the 7-feature ensemble against a multi-detector ensemble combining 4 detectors (FastDetectGPT, Binoculars, LD-Score, Avg-LogProb) plus surface features (Table~\ref{tab:ensemble}).
The multi-detector ensemble achieves a marginal improvement of $+$2.1 pp (3-class) and $+$0.3 pp (binary) at 2.4$\times$ inference cost and 3.2$\times$ GPU memory.
The binary \fone{} gap (0.3 pp) is essentially negligible, confirming that the gen-1/gen-2+ boundary is an intrinsic bottleneck that additional features cannot overcome.

\begin{table}[h]
\centering
\caption{Multi-detector ensemble comparison.}
\label{tab:ensemble}
\begin{tabular}{lccc}
\toprule
\textbf{Model} & \textbf{3-class \fone{}} & \textbf{Binary \fone{}} & \textbf{Cost Ratio} \\
\midrule
7-feature ensemble & 0.774 & 0.705 & 1.0$\times$ \\
4-det + surface & 0.795 & 0.708 & 2.4$\times$ \\
\bottomrule
\end{tabular}
\end{table}

\section{Finer-Grained Ordinal Classification}
\label{app:ordinal}

We test 5-class ordinal classification (gen-0 through gen-4; Table~\ref{tab:ordinal}).
All methods---standard XGBoost, CORAL, and Frank-Hall---achieve only $\sim$0.50 \fone{}, a $-$26 pp drop from 3-class.
The bottleneck is the gen-1/gen-2 boundary, which concentrates 86\% of binary false negatives.
Gen-3 vs.\ gen-4 is nearly indistinguishable (\fone{} $\approx$ 0.35).
Ordinal methods provide no improvement over standard classification, confirming that the resolution ceiling is at 3 classes---a feature-space limitation, not a methodological one.

\begin{table}[h]
\centering
\caption{Ordinal classification results.}
\label{tab:ordinal}
\begin{tabular}{lccc}
\toprule
\textbf{Method} & \textbf{3-class \fone{}} & \textbf{5-class \fone{}} & $\Delta$ \\
\midrule
XGBoost (standard) & 0.762 & 0.500 & $-0.262$ \\
CORAL & 0.758 & 0.490 & $-0.268$ \\
Frank-Hall & 0.760 & 0.500 & $-0.260$ \\
\bottomrule
\end{tabular}
\end{table}

\section{Temperature Sensitivity}
\label{app:temp}

Table~\ref{tab:temp} shows that DDR is reliable within $t \in [0.5, 0.7]$ but degrades substantially at $t = 0.9$.
The entropy\_mean Spearman $\rhosp$ drops from $-0.600$ (at $t = 0.7$) to $-0.302$ (at $t = 0.9$):
high-temperature sampling noise drowns out the deterministic depth-wise drift component.

\begin{table}[h]
\centering
\caption{Temperature sensitivity.}
\label{tab:temp}
\begin{tabular}{lccc}
\toprule
\textbf{Temperature} & \textbf{Within-CV Bin.\ \fone{}} & \textbf{LOGO Bin.\ \fone{}} & \textbf{entropy\_mean $\rhosp$} \\
\midrule
0.5 & 0.843 & 0.814 & $-0.584$ \\
0.7 (default) & 0.947 & 0.694 & $-0.600$ \\
0.9 & 0.582 & 0.480 & $-0.302$ \\
\bottomrule
\end{tabular}
\end{table}

\section{Human Edit Robustness}
\label{app:edit}

Table~\ref{tab:edit} shows graceful degradation under simulated human editing.
At 5\% random word replacement, binary \fone{} drops only 4.6 pp.
LLM rephrasing (average 1.69 sentences replaced, 50.7\% edit distance) produces a comparable drop ($-$4.5 pp),
suggesting that LLM editing does not introduce out-of-distribution perturbations in feature space.
Even at 20\% replacement, performance remains above random (\fone{} = 0.591).

\begin{table}[h]
\centering
\caption{Robustness to text editing.}
\label{tab:edit}
\begin{tabular}{lcc}
\toprule
\textbf{Condition} & \textbf{Binary \fone{}} & $\Delta$ \textbf{vs baseline} \\
\midrule
No editing (baseline) & 0.695 & --- \\
5\% random replacement & 0.650 & $-0.046$ \\
10\% random replacement & 0.614 & $-0.081$ \\
20\% random replacement & 0.591 & $-0.104$ \\
LLM rephrase (avg 1.69 sent) & 0.647 & $-0.045$ \\
\bottomrule
\end{tabular}
\end{table}

\section{LD-Score Augmentation}
\label{app:ld}

Adding LD-Score as an 8th feature yields a marginal improvement of $+$0.5 pp (Table~\ref{tab:ld}).
LD-Score alone achieves LOGO \fone{} = 0.721 (Spearman $\rhosp = -0.794$ with depth),
confirming character-level information partially captures depth signal but is highly correlated with existing reference features.
We retain the 7-feature set as the primary model for simplicity.

\begin{table}[h]
\centering
\caption{LD-Score augmentation.}
\label{tab:ld}
\begin{tabular}{lcc}
\toprule
\textbf{Feature Set} & \textbf{LOGO \fone{}} & $\Delta$ \\
\midrule
7-feature (default) & 0.798 & --- \\
7-feature + LD-Score & 0.803 & $+0.005$ \\
LD-Score only & 0.721 & $-0.077$ \\
\bottomrule
\end{tabular}
\end{table}

\section{End-to-End QLoRA Classifier}
\label{app:qlora}

Table~\ref{tab:qlora} compares QLoRA fine-tuned models against the 7-feature ensemble.
QLoRA models achieve 3-class \fone{} of 0.859--0.907, surpassing the 7-feature ensemble by 10--15 pp.
However, the 5-class ceiling persists (\fone{} $\approx$ 0.60), confirming that the resolution limit is task-intrinsic rather than feature-engineering-specific.
QLoRA inference costs $\sim$15$\times$ more than the 7-feature pipeline (which requires only a single 410M forward pass),
positioning the 7-feature ensemble on efficiency advantage.

\begin{table}[h]
\centering
\caption{End-to-end QLoRA vs.\ 7-feature ensemble.}
\label{tab:qlora}
\begin{tabular}{llccc}
\toprule
\textbf{Model} & \textbf{Type} & \textbf{3-class \fone{}} & \textbf{Binary \fone{}} & \textbf{5-class \fone{}} \\
\midrule
Qwen3-1.7B (instruct) & QLoRA & 0.907 & 0.867 & 0.551 \\
OLMo-1B (instruct) & QLoRA & 0.888 & 0.836 & 0.598 \\
Qwen3-1.7B (base) & QLoRA & 0.904 & 0.860 & 0.554 \\
OLMo-1B (base) & QLoRA & 0.859 & 0.793 & 0.540 \\
\midrule
7-feature ensemble & XGBoost & 0.762 & 0.694 & 0.500 \\
\bottomrule
\end{tabular}
\end{table}

\section{Text Length Confound Control}
\label{app:length}

Text length shows weak correlation with depth (Spearman $\rhosp \in [-0.066, +0.134]$), ruling it out as a primary confound.
Length-normalized features (TTR, hapax computed on fixed-length windows) preserve and even strengthen DDR trends.
Length-matched evaluation (Table~\ref{tab:length}) shows a $-$4.6 pp drop (0.770 $\to$ 0.724), but a random subsample control drops by the same amount,
confirming this reflects deconfounding rather than information loss.

\begin{table}[h]
\centering
\caption{Length confound analysis.}
\label{tab:length}
\begin{tabular}{lcc}
\toprule
\textbf{Condition} & \textbf{Binary \fone{}} & $\Delta$ \textbf{vs unmatched} \\
\midrule
Unmatched (full) & 0.770 & --- \\
Length-matched & 0.724 & $-0.046$ \\
Random subsample control & 0.724 & $-0.046$ \\
\bottomrule
\end{tabular}
\end{table}

\section{9-Generator Scale Evaluation}
\label{app:9gen}

Extending to 9 generators (adding Gemma-2-2B, Llama-3.2-3B, Mistral-7B-v0.3; 42,015 total records),
LOGO 3-class \fone{} is 0.751 ($-$1.9 pp vs.\ 8-gen; Table~\ref{tab:9gen}).
Gemma-2-2B is a notable outlier (3-class 0.628, binary 0.440) due to SentencePiece tokenizer mismatch with the BPE-based reference model.
Mistral-7B (3-class 0.605) reflects 7B-scale signal compression.
Excluding the Gemma anomaly, all generators exceed \fone{} = 0.60.

\begin{table}[h]
\centering
\caption{9-generator per-generator LOGO \fone{}.}
\label{tab:9gen}
\small
\begin{tabular}{lccc}
\toprule
\textbf{Generator} & \textbf{3-class \fone{}} & \textbf{Binary \fone{}} & \textbf{Family} \\
\midrule
Pythia-1B & 0.857 & 0.724 & Pythia \\
Qwen3-0.6B & 0.819 & 0.710 & Qwen \\
Pythia-2.8B & 0.812 & 0.700 & Pythia \\
Qwen3-1.7B & 0.811 & 0.698 & Qwen \\
OLMo-2-1B & 0.773 & 0.663 & OLMo \\
Llama-3.2-3B & 0.759 & 0.680 & Llama \\
Qwen3-8B & 0.696 & 0.647 & Qwen \\
Gemma-2-2B & 0.628 & 0.440 & Gemma \\
Mistral-7B-v0.3 & 0.605 & 0.633 & Mistral \\
\bottomrule
\end{tabular}
\end{table}

\section{Prompt Type and Chain Type Comparison}
\label{app:prompt}

\paragraph{Prompt types.}
Table~\ref{tab:prompt} compares three prompt types on Pythia-2.8B.
Continuation shows the strongest within-type signal (CV \fone{} = 0.781) but lowest depth retention (83.1\%),
while elaboration has the highest retention (98.6\%) but weakest signal (0.730).
This is consistent with the DDR location-shift interpretation: stronger semantic constraints narrow the output distribution, reducing cross-depth variation.
Leave-one-prompt-type-out (LOPTO) average \fone{} is 0.753, only 0.4 pp below within-type average, indicating high cross-prompt-type transferability.

\begin{table}[h]
\centering
\caption{Prompt-type comparison (Pythia-2.8B).}
\label{tab:prompt}
\begin{tabular}{lcccc}
\toprule
\textbf{Prompt Type} & \textbf{CV \fone{}} & \textbf{LOPTO \fone{}} & \textbf{Retention} & \textbf{hapax $\rhosp$} \\
\midrule
Continuation & 0.781 & 0.762 & 0.831 & $-0.74$ \\
Inspired rewriting & 0.737 & 0.733 & 0.939 & $-0.80$ \\
Elaboration & 0.730 & 0.762 & 0.986 & $-0.83$ \\
\bottomrule
\end{tabular}
\end{table}

\paragraph{Paraphrastic vs.\ non-paraphrastic chains.}
Table~\ref{tab:para} compares continuation and paraphrasing chains.
Paraphrasing chains show substantially weaker depth signal (CV \fone{} = 0.62 vs.\ 0.781, $-$16.1 pp)
because semantic-preservation constraints induce attractor convergence by depth 2,
suppressing the distributional drift that DDR exploits.
This validates the task formulation: depth estimation is meaningful specifically for non-paraphrastic chains.

\begin{table}[h]
\centering
\caption{Paraphrastic vs.\ non-paraphrastic chains.}
\label{tab:para}
\begin{tabular}{lcccc}
\toprule
\textbf{Chain Type} & \textbf{CV \fone{}} & \textbf{hapax $\rhosp$} & \textbf{entropy\_mean $\rhosp$} & \textbf{Retention} \\
\midrule
Continuation & 0.781 & $-0.74$ & $-0.56$ & 0.831 \\
Paraphrasing & 0.620 & $-0.25$ & $-0.18$ & 0.970 \\
\midrule
$\Delta$ & $-0.161$ & $+0.49$ & $+0.38$ & $+0.139$ \\
\bottomrule
\end{tabular}
\end{table}

\section{TOST Bootstrap Stationarity Test}
\label{app:tost}

We use Two One-Sided Tests (TOST) with 1,000 bootstrap iterations and equivalence margin $d = 0.2$
to formally test feature stationarity (Table~\ref{tab:tost}).
Three features (entropy\_max, entropy\_min, tail\_decay) are confirmed stationary (100\% rejection rate, median $|d| < 0.1$).
entropy\_std is rejected (median $d = 0.458$, 0\% rejection rate), classifying it as an intermediate ``weak drifter''
despite its low Spearman $\rhosp$ ($-0.12$).
This refines the DDR taxonomy: 3 strong drifters + 3 true stationary + 1 intermediate.

\begin{table}[h]
\centering
\caption{TOST bootstrap stationarity test.}
\label{tab:tost}
\begin{tabular}{lcccc}
\toprule
\textbf{Feature} & \textbf{Median $d$} & \textbf{TOST $p$} & \textbf{Rej.\ Rate} & \textbf{Status} \\
\midrule
entropy\_max & 0.062 & 0.000 & 1.000 & Stationary \checkmark \\
entropy\_min & $-0.016$ & 0.000 & 1.000 & Stationary \checkmark \\
tail\_decay & 0.088 & 0.000 & 1.000 & Stationary \checkmark \\
entropy\_std & 0.458 & 1.000 & 0.000 & Intermediate \\
\bottomrule
\end{tabular}
\end{table}

\section{Independent Seed Validation}
\label{app:seed}

To rule out seed-sharing leakage, we evaluate with independent seed sets per generator (Table~\ref{tab:seed}).
The average binary \fone{} (0.680) differs from shared-seed LOGO (0.694) by only 1.4 pp,
within bootstrap confidence intervals, confirming that shared seeds do not inflate performance.

\begin{table}[h]
\centering
\caption{Independent vs.\ shared seed evaluation.}
\label{tab:seed}
\begin{tabular}{lcc}
\toprule
\textbf{Condition} & \textbf{Binary \fone{}} & $\Delta$ \\
\midrule
Shared seed LOGO & 0.694 & --- \\
Independent seed (Pythia$\to$OLMo) & 0.695 & $+0.001$ \\
Independent seed (OLMo$\to$Pythia) & 0.666 & $-0.028$ \\
Independent seed avg & 0.680 & $-0.014$ \\
\bottomrule
\end{tabular}
\end{table}

\section{Topic Drift Control}
\label{app:topic}

Semantic similarity (cosine) increases monotonically with depth (0.622 $\to$ 0.683 $\to$ 0.801 $\to$ 0.859; Table~\ref{tab:topic}),
indicating that deeper generations converge semantically.
However, DDR feature--depth correlations remain significant after controlling for semantic similarity,
and shape features show near-zero correlation with similarity ($|\rhosp| < 0.17$),
confirming that DDR is not merely a proxy for topic drift.

\begin{table}[h]
\centering
\caption{Semantic similarity by depth transition.}
\label{tab:topic}
\begin{tabular}{lcc}
\toprule
\textbf{Transition} & \textbf{Mean Similarity} & \textbf{Std} \\
\midrule
depth 0 $\to$ 1 & 0.622 & 0.165 \\
depth 1 $\to$ 2 & 0.683 & 0.184 \\
depth 2 $\to$ 3 & 0.801 & 0.201 \\
depth 3 $\to$ 4 & 0.859 & 0.190 \\
\bottomrule
\end{tabular}
\end{table}

\end{document}
